\clearpage\thispagestyle{empty}
\pdfbookmark[0]{Sources for Wire Spacings (November 1823)}{wire1823-sources}
\begin{center}\textbf{Sources for Wire Spacings (November 1823)}\end{center}
\noindent This AI translation follows Gauss’s \textit{Fadenintervalle} in \href{https://gdz.sub.uni-goettingen.de/id/PPN236007467}{\textit{Werke}, vol.~VI} (Göttingen, 1874), pp.~445--448. The paper first appeared as \textit{Neue Methode, die gegenseitigen Abstände der Fäden in Meridian-Fernröhren zu bestimmen} in H.~C. Schumacher’s \textit{Astronomische Nachrichten}, vol.~II, no.~43 (November 1823), cols.~\href{https://api.digitale-sammlungen.de/iiif/image/v2/bsb10049427_00218/full/1800,/0/default.jpg}{371--372}, \href{https://api.digitale-sammlungen.de/iiif/image/v2/bsb10049427_00219/full/1800,/0/default.jpg}{373--374} and \href{https://api.digitale-sammlungen.de/iiif/image/v2/bsb10049427_00220/full/1800,/0/default.jpg}{375--376}. The links lead to the Bavarian State Library’s copy of the volume, published in Altona in 1824.
\par\smallskip\noindent The first printing credits Gauss as \textit{Hofrath und Ritter} and carries his signature. The collected edition has the heading \textit{Fadenintervalle} and spells the lens maker’s name \textit{Frauenhofer}, where the first printing has \textit{Fraunhofer}. On 23 October, the second true distance for wire I is printed as 10 24,06 in col.~374. The translation retains 10 24.86 as transcribed from the damaged digit in the collected edition. The table also follows the collected edition’s placement of unit marks. The title \textit{Theoria Combinationis Observationum} is Gauss’s own abbreviated reference. Translation and typesetting produced by AI.
\clearpage
